\begin{lemma}
Let \(I\subseteq[n]\) satisfy
\[
|I|=k=\noderef{TwoBiteNaturalIndependenceScale}(1+\varepsilon,n),
\]
and assume \(0\le\varepsilon_1\le1\),
\(\noderef{ParameterHierarchyEventualComparisons}
(\varepsilon,\varepsilon_1,\varepsilon_2,n_0)\), \(n>n_0\), and
\[
m=\noderef{TwoBiteNaturalReducedVertexCount}(n),\qquad
p=\noderef{TwoBiteEdgeProbability}\!\left(\frac12,n\right).
\]
Suppose that every sample \(\omega\) satisfying both
\[
\noderef{TwoBiteRegularityEvent}(k,\varepsilon,\varepsilon_1,\varepsilon_2,
\omega)
\]
and
\[
\noderef{TwoBiteProjectionSizeEvent}(I,k,\ell_R,\ell_B,\omega)
\]
also satisfies
\[
\noderef{ProjectionOpenPairFunction}(\ell_R,\ell_B,k,p,n)
  -10\varepsilon_1k^2
\le
\left|\noderef{TwoBiteStagedOpenPairs}(\omega,\varepsilon,I)\right|.
\tag{1}
\]
Then, under \(\noderef{TwoBiteConditionalProbability}(n,m,p;-|-)\) and
conditioning on the projection-size cell alone, the probability that \(I\) is
independent in the final two-bite graph and that the regularity event holds is
at most
\[
\exp\!\left(
-p\bigl(\noderef{ProjectionOpenPairFunction}(\ell_R,\ell_B,k,p,n)
       -\varepsilon^3k^2\bigr)\right).
\]
\end{lemma}
\begin{proof}
Let
\[
S(\omega)=
\noderef{TwoBiteProjectionSizeEvent}(I,k,\ell_R,\ell_B,\omega)
\]
be the projection-size conditioning event, and let
\[
A(\omega)=
\bigl(I\text{ is independent in }\noderef{TwoBiteFinalGraph}(\omega)\bigr)
\wedge
\noderef{TwoBiteRegularityEvent}(k,\varepsilon,\varepsilon_1,\varepsilon_2,
\omega)
\]
be the numerator event in the Lean statement.  Thus the left hand side to be
bounded is exactly
\[
\noderef{TwoBiteConditionalProbability}(n,m,p; A\mid S).
\]
For readability also write
\[
R(\omega)=
\noderef{TwoBiteRegularityEvent}(k,\varepsilon,\varepsilon_1,\varepsilon_2,\omega)
\]
for the regularity conjunct of \(A\).

Set
\[
f=
\noderef{ProjectionOpenPairFunction}
  ((\ell_R:\mathbb R),(\ell_B:\mathbb R),(k:\mathbb R),p,(n:\mathbb R)).
\tag{2}
\]
We instantiate \(\noderef{FixedSetConditionalExposureCellBound}\) with this
same fixed set \(I\), the same integers \(n,m,k,\ell_R,\ell_B\), the same real
parameters \(p,\varepsilon,\varepsilon_1,\varepsilon_2\), and the same
\(n_0\).  Its scalar and hierarchy hypotheses are exactly the present
hypotheses:
\[
0\le\varepsilon_1,\quad \varepsilon_1\le1,\quad
\noderef{ParameterHierarchyEventualComparisons}
(\varepsilon,\varepsilon_1,\varepsilon_2,n_0),\quad n>n_0,
\]
\[
m=\noderef{TwoBiteNaturalReducedVertexCount}(n),\qquad
p=\noderef{TwoBiteEdgeProbability}\!\left(\frac12,n\right),
\]
\[
|I|=k,\qquad
k=\noderef{TwoBiteNaturalIndependenceScale}(1+\varepsilon,n).
\]
The remaining hypothesis of the child theorem is the pointwise implication
\[
\forall\omega,\quad
R(\omega)\to S(\omega)\to
f-10\varepsilon_1(k:\mathbb R)^2
\le
\left|\noderef{TwoBiteStagedOpenPairs}(\omega,\varepsilon,I)\right|.
\tag{3}
\]
This is exactly assumption (1), after replacing the displayed
\(\noderef{ProjectionOpenPairFunction}\) by the abbreviation \(f\) from (2).
The order of hypotheses is also the same as in
\(\noderef{FixedSetConditionalExposureCellBound}\): first the regularity
event \(R(\omega)\), then the projection-size event \(S(\omega)\), and then
the lower bound for
\(\left|\noderef{TwoBiteStagedOpenPairs}(\omega,\varepsilon,I)\right|\).
Thus the pointwise staged-open lower-bound hypothesis is not strengthened or
weakened; it is the child theorem's hypothesis with the same function
\[
f=\noderef{ProjectionOpenPairFunction}
  ((\ell_R:\mathbb R),(\ell_B:\mathbb R),(k:\mathbb R),p,(n:\mathbb R)).
\]
No conditioning event is changed in this step: the regularity hypothesis is
used only as the first premise \(R(\omega)\) of the pointwise implication,
and the denominator event remains the projection-size predicate \(S\).

The conclusion of \(\noderef{FixedSetConditionalExposureCellBound}\), after
this instantiation, is therefore the bound for exactly the same conditional
probability \(\noderef{TwoBiteConditionalProbability}(n,m,p;A\mid S)\):
\[
\noderef{TwoBiteConditionalProbability}(n,m,p;A\mid S)
\le
\exp\!\left(
-p f
+8\sqrt{\varepsilon_1}\,p(k:\mathbb R)^2
+10\varepsilon_1p(k:\mathbb R)^2
+4\log (k:\mathbb R)\right).
\tag{4}
\]
The numerator predicate here is word-for-word the statement's numerator:
\[
\omega\mapsto
\bigl((\noderef{TwoBiteFinalGraph}(\omega))\text{ has }I\text{ independent}\bigr)
\wedge
\noderef{TwoBiteRegularityEvent}(k,\varepsilon,\varepsilon_1,\varepsilon_2,
\omega).
\]
It is not the projection-size event and it is not regularity alone.  The
denominator predicate here is word-for-word the statement's denominator:
\[
\omega\mapsto
\noderef{TwoBiteProjectionSizeEvent}(I,k,\ell_R,\ell_B,\omega).
\]
It contains no regularity conjunct.  Thus (4) already bounds the desired
conditional probability, but with the larger exponent containing the explicit
error terms.

It remains only to compare exponents.  Let
\[
\Delta=
8\sqrt{\varepsilon_1}\,p(k:\mathbb R)^2
+10\varepsilon_1p(k:\mathbb R)^2
+4\log (k:\mathbb R).
\]
The eventual-comparisons package is evaluated at the present \(n>n_0\).  In
that definition, the rounded integer scale and edge probability are
\[
K=\noderef{TwoBiteNaturalIndependenceScale}(1+\varepsilon,n),
\]
and
\[
\frac12\sqrt{\frac{\log n}{n}}
=\noderef{TwoBiteEdgeProbability}\!\left(\frac12,n\right).
\]
The numerical conjunct of
\(\noderef{ParameterHierarchyEventualComparisons}\) needed here is the
displayed comparison which absorbs
\[
8\sqrt{\varepsilon_1}\,pk^2+10\varepsilon_1pk^2+4\log k,
\]
and, applied to this \(n\), reads
\[
8\sqrt{\varepsilon_1}\,
\noderef{TwoBiteEdgeProbability}\!\left(\frac12,n\right)\,(K:\mathbb R)^2
+10\varepsilon_1
\noderef{TwoBiteEdgeProbability}\!\left(\frac12,n\right)\,(K:\mathbb R)^2
+4\log (K:\mathbb R)
\le
\varepsilon^3
\noderef{TwoBiteEdgeProbability}\!\left(\frac12,n\right)\,(K:\mathbb R)^2.
\tag{5a}
\]
The present hypotheses
\[
k=K,\qquad
p=\noderef{TwoBiteEdgeProbability}\!\left(\frac12,n\right)
\]
rewrite (5a), including the real casts of \(k\), into
\[
\Delta\le \varepsilon^3p(k:\mathbb R)^2. \tag{5}
\]
Adding the same term \(-pf\) to both sides of (5) gives
\[
 -pf+\Delta
\le
-pf+\varepsilon^3p(k:\mathbb R)^2. \tag{6}
\]
The right side of (6) is the target exponent.  Indeed, by commutativity and
associativity of multiplication in \(\mathbb R\),
\[
\varepsilon^3p(k:\mathbb R)^2
=p\bigl(\varepsilon^3(k:\mathbb R)^2\bigr),
\]
and by distributivity,
\[
-pf+\varepsilon^3p(k:\mathbb R)^2
=-p\bigl(f-\varepsilon^3(k:\mathbb R)^2\bigr). \tag{7}
\]
Combining (6) with this algebraic identity, we have
\[
-pf+\Delta
\le
-p\bigl(f-\varepsilon^3(k:\mathbb R)^2\bigr).
\tag{8}
\]
The real exponential is monotone increasing; in Lean this is the content of
\(\texttt{Real.exp\_le\_exp}\), which states
\(\exp(x)\le\exp(y)\) if and only if \(x\le y\).  Applying this monotonicity
to (8) gives
\[
\exp(-pf+\Delta)
\le
\exp\!\left(-p\bigl(f-\varepsilon^3(k:\mathbb R)^2\bigr)\right).
\tag{9}
\]
Combining the child bound (4) with (9) yields
\[
\noderef{TwoBiteConditionalProbability}(n,m,p;A\mid S)
\le
\exp\!\left(-p\bigl(f-\varepsilon^3(k:\mathbb R)^2\bigr)\right).
\tag{10}
\]
Finally substitute the definition (2) of \(f\).  This is exactly
\[
\exp\!\left(
-(p\cdot(
\noderef{ProjectionOpenPairFunction}
  ((\ell_R:\mathbb R),(\ell_B:\mathbb R),(k:\mathbb R),p,(n:\mathbb R))
 -\varepsilon^3(k:\mathbb R)^2))\right),
\]
matching the stated bound.
This is precisely the claimed conditional probability bound.
\end{proof}
